\section{The non-normal fibre and the CDP obstruction}

Campana--Demailly--Peternell prove that a compact complex threefold with
$H^1(X;\mathbb Z)=H^2(X;\mathbb Z)=0$, algebraic dimension one, and
holomorphic algebraic reduction has $c_3(X)\le0$
\cite[Theorem~2.2, p.~680]{CDP2020}.  Their corollary gives $a(X)=0$ for
a compact complex threefold homeomorphic to $S^6$
\cite[Corollary~2.3, p.~680]{CDP2020}.  The candidate asserts all those
hypotheses together with $c_3(X)=2$ and $a(X)=1$.  This section separates:
\begin{enumerate}[label=(\roman*)]
\item validity of one printed inference on p.~692 of \cite{CDP2020};
\item validity of the published theorem by any other argument;
\item validity of the candidate's proposed non-normal-fibre mechanism; and
\item truth of the candidate's upstream construction and topology.
\end{enumerate}

\subsection{The exact p.~692 inference}

Let $S$ be the reduction of a fibre in a smooth threefold $\mathcal X$,
and let $A$ denote the line-bundle twist after duality.  The two relevant
sequences are
\[
0\longrightarrow N^*_{S/\mathcal X}\otimes A
\longrightarrow\Omega^1_{\mathcal X}|_S\otimes A
\xrightarrow{\rho}\Omega^1_S\otimes A\longrightarrow0
\tag{5.1}
\]
and
\[
0\longrightarrow\mathcal T(\Omega^1_S)\otimes A
\longrightarrow\Omega^1_S\otimes A
\longrightarrow\widetilde\Omega^1_S\otimes A\longrightarrow0,
\tag{5.2}
\]
where $\mathcal T$ denotes sections supported on a nowhere dense analytic
set.  The opening reduction of \cite[\S7, p.~692]{CDP2020} treats the
vanishing of the global sections of the left end of (5.1) and the right end
of (5.2) as sufficient for the vanishing of the middle term of (5.1).
That inference leaves untreated a section whose image under $\rho$ is a
nonzero torsion section.  The omitted sheaf can be identified exactly.  Put
\[
 \mathscr T_S:=\ker\bigl(\Omega_S^1\to
 j_*\Omega^1_{S_{\rm reg}}\bigr),
 \qquad \widetilde\Omega_S^1=\Omega_S^1/\mathscr T_S,
\]
and
\[
 \mathscr K_S:=
 \ker\bigl(\Omega_X^1|_S\longrightarrow\widetilde\Omega_S^1\bigr).
\]

\begin{proposition}[the missing cotangent filtration]
\label{prop:missing-cotangent-filtration}
Let $S$ be a reduced effective Cartier divisor in a smooth complex
threefold $X$.  There are natural short exact sequences
\[
0\longrightarrow N^*_{S/X}\longrightarrow\mathscr K_S
\longrightarrow\mathscr T_S\longrightarrow0,
\]
\[
0\longrightarrow\mathscr K_S\longrightarrow\Omega_X^1|_S
\longrightarrow\widetilde\Omega_S^1\longrightarrow0.
\]
\end{proposition}

\begin{proof}
The conormal sequence is right exact.  It is also left exact here.  Locally
write $S=(g=0)$ with $g$ reduced.  If $h\,dg|_S=0$, then $h$ vanishes at the
generic smooth point of every irreducible component, because $dg$ is nonzero
there.  The ring $\mathcal O_S$ is reduced, so the intersection of its
minimal primes is zero; hence $h=0$.  Thus
\[
0\to N^*_{S/X}\to\Omega_X^1|_S\to\Omega_S^1\to0
\]
is exact.  By definition, $\mathscr K_S$ is the inverse image of
$\mathscr T_S$ under the last surjection.  Restriction to that inverse image
is onto $\mathscr T_S$ and has kernel $N^*_{S/X}$, proving the first
sequence.  Composing with the quotient by $\mathscr T_S$ proves the second.
\end{proof}

\begin{proposition}[the exact global-section criterion]
\label{prop:missing-global-section-criterion}
Let $A$ be a line bundle on $S$.  If
\[
H^0(S,\widetilde\Omega_S^1\otimes A)=0,
\]
then
\[
H^0(S,\Omega_X^1|_S\otimes A)
=H^0(S,\mathscr K_S\otimes A).
\]
If also $H^0(S,N^*_{S/X}\otimes A)=0$, then this group is canonically
\[
\ker\!\left[
H^0(S,\mathscr T_S\otimes A)
\xrightarrow{\delta_A}H^1(S,N^*_{S/X}\otimes A)
\right].
\]
Hence the two vanishings printed in the p.~692 reduction imply the desired
ambient vanishing if and only if $\delta_A$ is injective.
\end{proposition}

\begin{proof}
Tensoring the second sequence of
\cref{prop:missing-cotangent-filtration} by the locally free sheaf $A$
preserves exactness; its global-section sequence gives the first equality.
Tensor the first sequence and take global sections.  If the conormal
$H^0$ vanishes, exactness identifies the middle group with the displayed
kernel.  No condition on the torsion-free quotient controls that kernel.
\end{proof}

\subsection{Local algebra of a reduced normal-crossings fibre}

Let $g:\mathcal X\to\Delta$ be a proper holomorphic map from a smooth complex
threefold to a disc, and put $S=g^{-1}(0)$.  Assume $S$ is reduced and has
normal crossings: at every point there are coordinates $(x,y,z)$ in which
\[
g=ux,\qquad g=uxy,\qquad\text{or}\qquad g=uxyz
\tag{5.3}
\]
for a unit $u$.  Let $\eta:\widetilde S\to S$ be the normalization,
$D=S_{\mathrm{sing,red}}$, and $\widetilde D=\eta^{-1}(D)_{\mathrm{red}}$.

\begin{lemma}[the conormal map is injective]\label{lem:conormal-injective}
\Verified\ In the setting above, $N^*_{S/\mathcal X}\simeq\mathcal O_S$,
and its generator maps to $dg|_S$.  Multiplication by $dg|_S$ is injective,
so
\[
0\longrightarrow\mathcal O_S
\xrightarrow{\,h\mapsto h\,dg\,}
\Omega^1_{\mathcal X}|_S\longrightarrow\Omega^1_S\longrightarrow0
\tag{5.4}
\]
is exact.
\end{lemma}

\begin{proof}
The fibre ideal is the principal ideal $(g)$ and is reduced, so
$(g)/(g^2)\simeq\mathcal O_S$.  Multiplication by a unit does not affect
injectivity, and the term $f\,du$ vanishes after restriction to $f=0$.
It therefore suffices to use $f=x,xy,xyz$.

For $f=xy$, if a lift $\widetilde h$ satisfies
$\widetilde h(y\,dx+x\,dy)=0$ modulo $(xy)$, then
$\widetilde h y,\widetilde h x\in(xy)$.  Hence
$\widetilde h\in(x)\cap(y)=(xy)$ in the regular local ring, so $h=0$ on
$S$.  For $f=xyz$, the three coefficient conditions say
$\widetilde h yz,\widetilde h xz,\widetilde h xy\in(xyz)$; hence
$\widetilde h\in(x)\cap(y)\cap(z)=(xyz)$.  The case $f=x$ is immediate.
Thus the conormal map is injective in every local model, proving (5.4).
\end{proof}

\begin{lemma}[normalization--conductor exact sequence]
\label{lem:conductor-square}
\Verified\ The sequence
\[
0\longrightarrow\mathcal O_S
\xrightarrow{\alpha}
\eta_*\mathcal O_{\widetilde S}\oplus\mathcal O_D
\xrightarrow{\beta}
\eta_*\mathcal O_{\widetilde D}
\longrightarrow0,
\tag{5.5}
\]
\[
\alpha(h)=(\eta^*h,h|_D),\qquad
\beta(\widetilde h,k)=
\widetilde h|_{\widetilde D}-\eta^*k,
\]
is exact.
\end{lemma}

\begin{proof}
The assertion is local.  Away from $D$, $\eta$ is an isomorphism.
At a double point $S=\{xy=0\}$, a member of $\ker\beta$ is
\[
\bigl((h_1(y,z),h_2(x,z)),k(z)\bigr),
\qquad
h_1(0,z)=h_2(0,z)=k(z).
\]
The function
\[
h=h_1(y,z)+h_2(x,z)-h_1(0,z)
\]
on $\{xy=0\}$ restricts to $h_1,h_2$ on the two branches and to $k$ on
their intersection.  This proves middle exactness in the double model.

At a triple point $S=\{xyz=0\}$, write the branch functions as
$h_1(y,z)$, $h_2(x,z)$, $h_3(x,y)$.  Membership in $\ker\beta$ is exactly
\[
h_1(0,z)=h_2(0,z),\quad
h_1(y,0)=h_3(0,y),\quad
h_2(x,0)=h_3(x,0).
\]
Then
\[
\begin{split}
h={}&h_1(y,z)+h_2(x,z)+h_3(x,y)\\
&-h_1(0,z)-h_1(y,0)-h_2(x,0)+h_1(0,0)
\end{split}
\]
restricts to the prescribed function on every branch.  Injectivity of
$\alpha$ follows because $S$ is reduced and $\eta$ is an isomorphism over a
dense open set.  Surjectivity of $\beta$ follows locally by lifting a
function on the branch conductor to the smooth branch and taking $k=0$.
This proves exactness in every local model.
\end{proof}

\begin{theorem}[conductor-supported ambient differential]
\label{prop:conductor}
\Verified\ In the setting (5.3), let $A$ be a line bundle on $S$ with
$\eta^*A\simeq\mathcal O_{\widetilde S}$, and choose a nowhere-vanishing
section $e$ of $\eta^*A$.  There is a unique nonzero
\[
s\in H^0(S,\Omega^1_{\mathcal X}|_S\otimes A)
\tag{5.6}
\]
such that
\[
\eta^*s=\eta^*(dg|_S)\otimes e,\qquad s|_D=0.
\tag{5.7}
\]
Its image $\rho(s)$ in $\Omega^1_S\otimes A$ is supported on $D$ and
therefore torsion; its image in
$\widetilde\Omega^1_S\otimes A$ is zero.  Moreover:
\begin{enumerate}[label=(\alph*)]
\item $\rho(s)=0$ if and only if $e=\eta^*\sigma$ for a nowhere-vanishing
$\sigma\in H^0(S,A)$;
\item if $S$ is connected and $A\not\simeq\mathcal O_S$, then
$H^0(S,A)=0$ and $s$ does not come from the conormal term in (5.1).
\end{enumerate}
\end{theorem}

\begin{proof}
In the double model, $d(xy)=y\,dx+x\,dy$ has coefficients in the reduced
ideal $(x,y)$ of $D$.  In the triple model,
\[
d(xyz)=yz\,dx+xz\,dy+xy\,dz
\]
has coefficients in $(xy,xz,yz)$, the reduced ideal of the union of the
three double curves.  Unit factors in (5.3) do not change this after
restriction to $S$.  Hence $\eta^*(dg|_S)$ restricts to zero on
$\widetilde D$.  Tensor (5.5) by the locally free sheaf
$E=\Omega^1_{\mathcal X}|_S\otimes A$.  The compatible pair
\[
\bigl(\eta^*(dg|_S)\otimes e,0\bigr)
\]
therefore descends uniquely to $s$.  On the smooth locus of $S$, $dg|_S$
is a nonzero conormal vector and $e$ is nowhere zero, so $s\ne0$.

On $S\setminus D$, the normalization is an isomorphism and
$s=dg|_S\otimes e$.  By (5.4) its intrinsic image vanishes there.  Thus
$\rho(s)$ is supported on $D$ and lies in the torsion subsheaf; its image
modulo torsion is zero.

If $\rho(s)=0$, exactness of (5.4) after tensoring by $A$ gives
$s=dg|_S\otimes\sigma$ for a section $\sigma$ of $A$.  Pulling back and
comparing with (5.7) on the dense open set where $dg\ne0$ gives
$\eta^*\sigma=e$ there, hence everywhere by holomorphic continuation.
Since $e$ never vanishes and $\eta$ is surjective, $\sigma$ never vanishes
and trivializes $A$.  The converse is immediate.

Finally suppose $S$ is connected and $\sigma\in H^0(S,A)$.  On each compact
connected component of $\widetilde S$, write
$\eta^*\sigma=c_i e$; the holomorphic function $c_i$ is constant.  If one
$c_i$ is zero and another is nonzero, connectedness of the finite union of
irreducible components of $S$ provides two such components meeting at a
point.  The value of $\sigma$ at that point would be both zero and nonzero,
a contradiction.  Hence either every $c_i=0$, so $\sigma=0$, or every
$c_i\ne0$, so $\sigma$ is nowhere zero and $A$ is trivial.  When $A$ is
nontrivial only the first case occurs.  Since
$N^*_{S/\mathcal X}\simeq\mathcal O_S$, the conormal global-section space is
$H^0(S,A)=0$, while $s\ne0$.
\end{proof}

\subsection{A proper counterexample to the printed inference}

\begin{theorem}[nodal-cubic test family]\label{prop:nodal-counterexample}
\Refuted\ The p.~692 implication described after (5.2) is false as a formal
sheaf-theoretic statement, even for a proper family with smooth total space
and an irreducible reduced normal-crossings central fibre.
\end{theorem}

\begin{proof}
Let $\Delta$ be a sufficiently small disc with coordinate $t$ and put
\[
\mathcal C=
\{([x:y:z],t)\in\mathbb P^2\times\Delta:
y^2z=x^2(x+z)+tz^3\}.
\tag{5.8}
\]
This is proper over $\Delta$.  Its total space is smooth.  On $z\ne0$, the
partial derivative of the defining equation with respect to $t$ is
$-z^3\ne0$.  On $z=0$, the equation forces $x=0$; then $y\ne0$ and the
$z$-derivative is $y^2\ne0$.  The central fibre $C_0$ is the irreducible
nodal cubic $y^2z=x^2(x+z)$, with node $[0:0:1]$.

Set
\[
\mathcal X=\mathcal C\times\mathbb P^1,\qquad
S=C_0\times\mathbb P^1,
\]
and let $g:\mathcal X\to\Delta$ be projection.  Then $\mathcal X$ is a
smooth threefold, $g$ is proper, and $S$ is an irreducible reduced
normal-crossings surface.  Its normalization is
$\widetilde S\simeq\mathbb P^1\times\mathbb P^1$.

Let $p,q$ be the two inverse images of the node in the normalization of
$C_0$.  Fix $\lambda\in\mathbb C^*\setminus\{1\}$.  Glue the fibres of
the trivial line bundle on the normalization at $p$ and $q$ by the scalar
$\lambda$.  This defines a line bundle
$A_\lambda\in\operatorname{Pic}^0(C_0)$ whose pullback is trivial.  It is
not trivial: a global section pulls back to a constant $c$ satisfying
$c=\lambda c$, so
\[
H^0(C_0,A_\lambda)=0.
\tag{5.9}
\]
Put $A=\operatorname{pr}_1^*A_\lambda$ on $S$.  Then
$\eta^*A\simeq\mathcal O_{\widetilde S}$ and $H^0(S,A)=0$.
Because $S$ is the full fibre,
$N^*_{S/\mathcal X}\simeq\mathcal O_S$, so the conormal-twist global
sections vanish.

The natural map
$\widetilde\Omega^1_S\to\eta_*\Omega^1_{\widetilde S}$ is injective:
its kernel is supported where normalization is not an isomorphism, while
the source is torsion-free by definition.  Projection formula and
$\eta^*A\simeq\mathcal O_{\widetilde S}$ therefore give
\[
H^0(S,\widetilde\Omega^1_S\otimes A)
\hookrightarrow
H^0(\widetilde S,\Omega^1_{\widetilde S})=0,
\tag{5.10}
\]
the last equality following from
$\widetilde S=\mathbb P^1\times\mathbb P^1$.
Nevertheless \cref{prop:conductor} produces
\[
0\ne s\in H^0(S,\Omega^1_{\mathcal X}|_S\otimes A).
\]
Thus both vanishings invoked in the printed reduction hold while the desired
ambient vanishing fails.  The family (5.8) does not satisfy the global
hypotheses of CDP Theorem~2.2, so this is a counterexample to the inference,
not to the theorem.
\end{proof}

\subsection{Cartier-divisor duality without a normality hypothesis}

\begin{theorem}[complete Cartier-fibre duality]
\label{thm:cartier-duality}
\Verified\ Let $X$ be a compact smooth complex threefold,
$i:W\hookrightarrow X$ an effective Cartier divisor, and $E$ a holomorphic
vector bundle.  Put $F=E|_W$ and
\[
\omega_W=(\omega_X\otimes\mathcal O_X(W))|_W.
\]
For $q=0,1,2$ there are natural perfect pairings
\[
H^q(W,F)^\vee
\simeq H^{2-q}(W,F^\vee\otimes\omega_W).
\tag{5.11}
\]
Moreover $H^q(W,F)=0$ for $q>2$.  No normality or reducedness hypothesis on
$W$ is used.
\end{theorem}

\begin{proof}
The Cartier divisor gives two short exact sequences
\[
0\to E(-W)\to E\to i_*F\to0
\tag{5.12}
\]
and, by adjunction,
\[
0\to E^\vee\otimes\omega_X
\to E^\vee\otimes\omega_X(W)
\to i_*(F^\vee\otimes\omega_W)\to0.
\tag{5.13}
\]
Smooth Serre duality on the compact threefold gives functorial isomorphisms
\[
H^r(X,G)^\vee\simeq H^{3-r}(X,G^\vee\otimes\omega_X)
\tag{5.14}
\]
for every vector bundle $G$
\cite[Theorem~9.37]{Lee2024}.  Naturality of the trace pairing says that the
dual of the map on cohomology induced by $E(-W)\to E$ is the map induced by
$E^\vee\otimes\omega_X\to E^\vee\otimes\omega_X(W)$.

Fix $q\in\{0,1,2\}$.  Dualize the relevant five terms of the long exact
sequence of (5.12) and use (5.14).  They become
\[
\begin{split}
H^{2-q}(X,E^\vee\otimes\omega_X)
&\longrightarrow H^{2-q}(X,E^\vee\otimes\omega_X(W))\\
&\longrightarrow H^q(W,F)^\vee\\
&\longrightarrow H^{3-q}(X,E^\vee\otimes\omega_X)
\longrightarrow H^{3-q}(X,E^\vee\otimes\omega_X(W)).
\end{split}
\tag{5.15}
\]
The long exact sequence of (5.13) is the same five-term sequence with its
middle group replaced by
$H^{2-q}(W,F^\vee\otimes\omega_W)$.  Naturality supplies a commutative
diagram whose four outer vertical arrows are identities.  The five lemma
(or exactness at the three middle terms) gives the unique middle isomorphism
(5.11).  Finite dimensionality makes the pairing perfect.

For $q=3$, the map
$H^3(X,E(-W))\to H^3(X,E)$ is dual under (5.14) to the injective map
\[
H^0(X,E^\vee\otimes\omega_X)
\hookrightarrow H^0(X,E^\vee\otimes\omega_X(W)).
\]
It is therefore surjective.  The long exact sequence of (5.12) gives
$H^3(W,F)=0$.  For $q>3$ both adjacent ambient cohomology groups vanish by
the dimension-three Dolbeault resolution, so (5.12) gives
$H^q(W,F)=0$.  This also agrees with analytic duality for the Gorenstein
hypersurface $W$ in Ramis--Ruget
\cite[Theorems~1--2 and Proposition~2]{RamisRuget1970}.
\end{proof}

\subsection{Flatness and top-degree analytic base change}

\begin{lemma}[a nonconstant map from a manifold to a curve is flat]
\label{lem:curve-flatness}
\Verified\ Let $f:X\to B$ be a nonconstant holomorphic map from a connected
complex manifold to a Riemann surface.  Then $f$ is flat.  Every holomorphic
vector bundle on $X$ is $f$-flat.
\end{lemma}

\begin{proof}
Fix $x\in X$, $c=f(x)$, and put
$R=\mathcal O_{B,c}$, $S=\mathcal O_{X,x}$.  The ring $R$ is a discrete
valuation ring and $S$ is a regular local domain.  The homomorphism
$R\to S$ is injective.  Indeed, if a nonzero germ $h\in R$ mapped to zero,
$h\circ f$ would vanish on a neighbourhood of $x$.  The zero set of a
nonzero one-variable holomorphic germ is isolated, so $f$ would be constant
on that neighbourhood and therefore, by the identity theorem on connected
$X$, constant everywhere.

Thus every nonzero element of $R$ acts on the domain $S$ as a nonzero
element and hence as a nonzerodivisor.  Therefore $S$ is torsion-free as an
$R$-module.  A module over a discrete valuation ring is flat if and only if
it is torsion-free \cite[Tag~0AUW]{StacksProject}; hence $S$ is $R$-flat.
This holds at every point, so $f$ is flat.  A vector bundle is locally a
finite free $\mathcal O_X$-module, and a finite direct sum of $R$-flat
modules is flat; hence it is $f$-flat.
\end{proof}

\begin{proposition}[the exact analytic base-change input]
\label{prop:analytic-base-change}
\Verified\ Let $f:X\to B$ be proper between complex spaces, let
$\mathcal F$ be coherent and $f$-flat, and fix $c\in B$.  If
\[
H^3(X_c,\mathcal F|_{X_c})=0,
\tag{5.16}
\]
then the canonical fibre map
\[
\beta_c^2:
(R^2f_*\mathcal F)_c\otimes_{\mathcal O_{B,c}}\mathbb C(c)
\longrightarrow H^2(X_c,\mathcal F|_{X_c})
\tag{5.17}
\]
is an isomorphism.
\end{proposition}

\begin{proof}
Apply B\u{a}nic\u{a}--St\u{a}n\u{a}\c{s}il\u{a},
Chapter~III, \S3, Corollary~3.5, with $q=3$
\cite[pp.~120--121]{BanicaStanasila1976}.  Its hypotheses are precisely
properness, coherence, and flatness.  Condition (5.16) first forces the
degree-three direct image to vanish near $c$.  The corollary then applies the
equivalent base-change conditions of its preceding Theorem~3.4 in degree
$q-1=2$ and states that the canonical map (5.17) is bijective.  This is the
top-degree argument used in \cite[Lemma~1.5]{CDP1998}; it does not require
the stronger, false assertion that all direct images are locally free.
\end{proof}

\subsection{Conditional theorem for the candidate fibre}

\begin{theorem}[the candidate conductor mechanism, with all hypotheses]
\label{thm:candidate-conductor}
\Conditional\ Assume:
\begin{enumerate}[label=(H\arabic*)]
\item $X$ is a connected compact smooth complex threefold and
$f:X\to B$ is a nonconstant holomorphic map to a smooth compact curve;
\item $W=f^{-1}(p_0)$ is the full scheme-theoretic fibre, is reduced, and
has the normal-crossings local models (5.3);
\item the normalization $\eta:\widetilde W\to W$ is the degree-six del
Pezzo surface $\mathrm{dP}_6$;
\item $H^2(X;\mathbb Z)=0$.
\end{enumerate}
Then, for every $L\in\operatorname{Pic}(X)$,
\[
H^2(W,(T_X\otimes L)|_W)\ne0
\quad\text{and}\quad
R^2f_*(T_X\otimes L)\ne0.
\tag{5.18}
\]
\end{theorem}

\begin{proof}
Fix $L$ and put
\[
\mathcal F_L=T_X\otimes L,\qquad
A_L=(L^\vee\otimes\omega_X)|_W.
\]
Hypothesis (H4) gives $c_1(L^\vee\otimes\omega_X)=0$.  Pullback therefore
gives $c_1(\eta^*A_L)=0$.  The surface $\mathrm{dP}_6$ is the blow-up of
$\mathbb P^2$ at three noncollinear points; the standard blow-up calculation
gives $H^1(\mathrm{dP}_6,\mathcal O)=0$
\cite[Chapter~V]{Hartshorne1977}.  In the exponential sequence, the kernel
of
\[
c_1:\operatorname{Pic}(\mathrm{dP}_6)\longrightarrow
H^2(\mathrm{dP}_6;\mathbb Z)
\]
is the image of $H^1(\mathrm{dP}_6,\mathcal O)$ and is therefore zero.
Consequently
\[
\eta^*A_L\simeq\mathcal O_{\widetilde W}.
\tag{5.19}
\]
\Cref{prop:conductor} now gives
\[
H^0(W,\Omega^1_X|_W\otimes A_L)\ne0.
\tag{5.20}
\]

Because $W=f^*(p_0)$, one has
$\mathcal O_X(W)=f^*\mathcal O_B(p_0)$, whose restriction to the fibre is
trivial.  Adjunction therefore gives $\omega_W\simeq\omega_X|_W$.
Apply \cref{thm:cartier-duality} to $E=\mathcal F_L$:
\[
\begin{split}
H^2(W,\mathcal F_L|_W)^\vee
&\simeq H^0(W,(\mathcal F_L|_W)^\vee\otimes\omega_W)\\
&=H^0(W,\Omega^1_X|_W\otimes A_L).
\end{split}
\]
The right side is nonzero by (5.20), proving the first assertion of (5.18).
The same duality theorem gives
$H^3(W,\mathcal F_L|_W)=0$.

Since $X$ is compact, $f$ is proper.  By
\cref{lem:curve-flatness}, $f$ is flat and the vector bundle
$\mathcal F_L$ is $f$-flat.  Thus every hypothesis of
\cref{prop:analytic-base-change} is satisfied at $p_0$, and
\[
(R^2f_*\mathcal F_L)_{p_0}\otimes\mathbb C(p_0)
\simeq H^2(W,\mathcal F_L|_W)\ne0.
\]
The stalk, and hence the coherent direct-image sheaf, is nonzero.  This is
the second assertion of (5.18).
\end{proof}

\begin{corollary}[application to the constructed threefold]
\label{cor:candidate-conductor-actual}
\Verified\ For the global construction of
\cref{thm:candidate-global-geometry}, the cusp hypotheses (H1)--(H3) and
the integral-topology hypothesis (H4) hold for the same threefold.  Hence
for every $L\in\operatorname{Pic}(X)$,
\[
H^2(W,(T_X\otimes L)|_W)\ne0,
\qquad R^2f_*(T_X\otimes L)\ne0.
\]
\end{corollary}

\begin{proof}
Hypotheses (H1)--(H3) are
\cref{thm:candidate-cusp-geometry,thm:candidate-global-geometry}, and (H4)
is \cref{thm:h2-vanishing}.  Apply
\cref{thm:candidate-conductor}.
\end{proof}

\subsection{A second printed issue: abelianization of split extensions}

\begin{lemma}[abelianization of a semidirect product]
\label{lem:semidirect-abelianization}
\Verified\ If
$1\to K\to G\to Q\to1$ is split, then
\[
G^{\mathrm{ab}}\simeq
Q^{\mathrm{ab}}\oplus(K^{\mathrm{ab}})_Q,
\tag{5.21}
\]
where the subscript denotes coinvariants for the conjugation action of $Q$.
\end{lemma}

\begin{proof}
Choose a splitting and write $G=K\rtimes Q$.  Let $N\triangleleft K$ be
generated by $[K,K]$ and by all elements
$qkq^{-1}k^{-1}$ with $q\in Q$, $k\in K$.  It is normal in $G$:
$[K,K]$ makes it stable under conjugation by $K$, while conjugation by
$Q$ permutes the displayed cross-commutators.  By construction
$K/N=(K^{\mathrm{ab}})_Q$, and $Q$ acts trivially on this quotient.  Hence
\[
G/N\simeq(K^{\mathrm{ab}})_Q\times Q.
\]
Since $N\subseteq[G,G]$, abelianizing the last display gives (5.21).
\end{proof}

The sentence in \cite[Lemma~3.2, pp.~84--85]{CDP1998} that passes from a
split fundamental-group extension to a direct sum with $H_1(F;\mathbb Z)$
is therefore valid only when monodromy acts trivially after abelianization.
A minimal counterexample retaining a free quotient group is
\[
G=\mathbb Z\rtimes_{-1}\mathbb Z.
\]
The extension splits, but
$G^{\mathrm{ab}}\simeq\mathbb Z\oplus\mathbb Z/2$, whereas the printed
direct-product formula gives $\mathbb Z^2$.
For the candidate's printed monodromy, \cref{thm:coinvariants} gives
\[
H^1(F;\mathbb Q)^{\pi_1(B^\circ)}=\mathbb Q\gamma,
\]
so the rational coinvariant rank is one, not
$b_1(F)=4$.

If $C^\circ$ is a sphere punctured at $s$ points, the correct rank formula is
\[
b_1(X^\circ)=s-1+
\dim_{\mathbb Q}H_1(F;\mathbb Q)_{\pi_1(C^\circ)}.
\tag{5.22}
\]
Thus the proof of CDP98 Theorem~3.1 reaches only the coinvariant rank unless
an additional theorem proves trivial monodromy.  CDP20 Lemma~4.2 repeats the
direct-sum formula, and CDP20 Corollary~5.2 imports CDP98 Theorem~3.1.
Repairing the local-freeness issue discussed there does not repair this
distinct monodromy error.  The torus branch of the CDP20 proof therefore has
an unresolved dependency in addition to the singular-fibre defect above.

\subsection{Exact present disposition}

\begin{enumerate}[label=\arabic*.]
\item \Refuted\ The p.~692 reduction in \cite{CDP2020} is not a formal
consequence of the two cotangent sequences.
\Cref{prop:nodal-counterexample} is a proper geometric counterexample to
that exact inference.
\item \Refuted\ The direct-product homology formula in
\cite[Lemma~3.2]{CDP1998} omits monodromy coinvariants.  This leaves the
torus branch imported by \cite{CDP2020} without its printed rank argument.
\item \Verified\ The normalization--conductor construction,
Cartier-fibre duality, flatness over the base curve, and top-degree analytic
base change are proved above under their stated hypotheses.
\item \Verified\ The period-family descent, toric and finite fillings,
global gluing, and integral analytic nearby-cycle/Leray bridges are proved in
\cref{sec:candidate-geometry,sec:integral-leray}.  Therefore the candidate's
hypotheses (H1)--(H4) hold and
$R^2f_*(T_X\otimes L)\ne0$ for every $L$ by
\cref{cor:candidate-conductor-actual}.
\item \Verified\ Those same construction and topology results, together
with \cref{thm:algebraic-dimension-audit}, hold for the constructed compact
complex threefold.  Hence
$H^1(X;\mathbb Z)=H^2(X;\mathbb Z)=0$, $a(X)=1$, $f$ is a holomorphic
algebraic reduction, and $c_3(X)=2$.  The resulting $X$ satisfies every
stated hypothesis of CDP Theorem~2.2 and contradicts its conclusion; it is
therefore a counterexample to that theorem as stated.  This theorem-level
conclusion uses the complete construction and topology proofs; it is not
inferred from the two printed proof defects alone.
\item \Open\ No independent repair of CDP Theorem~2.2 that bypasses both
refuted dependencies is recorded here.  Because the counterexample would be
major, independent human verification of every candidate proof above remains
essential; that social verification requirement is not disguised as a
mathematical hypothesis.
\end{enumerate}

\begin{test}[decisive theorem-level outcomes]\label{obl:cdp-comparison}
The current workbench records two exact defects in the printed CDP chain and
a completed route from the constructed threefold to a theorem-level
counterexample.  The remaining independent challenge is to falsify a named
step in that route or find an independent repair of CDP Theorem~2.2 that
proves one of the established candidate assertions impossible.  Merely
invoking the published theorem, without repairing its dependencies, does not
identify such a step; neither would the two local printed-proof defects by
themselves have decided the published theorem before the candidate was
proved.
\end{test}
